% ---------------------------------------------------------------------
% Companion note to paper1/main.tex, Secs. 6 and 8 (thm:dictionary,
% thm:channels, thm:kato, thm:stability, prop:directintegral) and to the
% 1D construction research/p1_thm_dictionary_feller_leg_proof.tex
% ("the 1D Lemma" below).
%
% QUESTION.  Does the escape-diffusion / Feller-boundary dictionary
% proved in 1D extend to the multidimensional / channel-based setting?
%
% ANSWER (proved here, honest labels throughout):
%   - separable and central cases: FULL per-channel generalization
%     (Theorem A, Corollary B), including the lateral-condition mirror
%     of deficiency (infinity,infinity);
%   - bounded channel mixing: dictionary survives (Theorem C), the
%     stochastic mirror of Kato stability (paper's thm:kato);
%   - genuine obstruction: UNBOUNDED CHANNEL TRANSFER.  The per-channel
%     integrals {T_n} provably do NOT determine the coupled
%     classification (Propositions E1, E2); the failure to explode is
%     "channel riding", and the failure to remain conservative is
%     ghost-specific (Theorem F: impossible when h_perp is bounded
%     below).  The channel tail n = infinity is a second Feller
%     boundary with its own one-sum criterion (Proposition G).
% VERDICT at the end of the file: PER-CHANNEL-ONLY (with the coupled
% classification reduced to a clean conjecture, and the quantum tail
% open exactly where the paper's Open Problem sits).
% ---------------------------------------------------------------------
\documentclass[11pt]{article}
\usepackage{amsmath,amssymb,amsthm}
\usepackage[margin=1.05in]{geometry}

\newtheorem{theorem}{Theorem}
\newtheorem{lemma}[theorem]{Lemma}
\newtheorem{proposition}[theorem]{Proposition}
\newtheorem{corollary}[theorem]{Corollary}
\newtheorem{conjecture}[theorem]{Conjecture}
\theoremstyle{remark}
\newtheorem{remark}[theorem]{Remark}

\newcommand{\dd}{\mathrm{d}}
\newcommand{\E}{\mathbb{E}}
\newcommand{\Prob}{\mathbb{P}}
\newcommand{\R}{\mathbb{R}}
\newcommand{\N}{\mathbb{N}}
\newcommand{\calE}{\mathcal{E}}
\newcommand{\ii}{\mathrm{i}}

\title{The escape dictionary beyond one dimension:\\
channel diffusions, switched escape processes,\\ and the exact
obstruction}
\author{Companion note to Theorem~6.1 and Secs.~8--9 of the paper\\
(builds on the 1D Lemma of
\texttt{p1\_thm\_dictionary\_feller\_leg\_proof.tex})}
\date{}

\begin{document}
\maketitle

\noindent\textbf{Status conventions.}  Every statement below is tagged
\emph{[proved]}, \emph{[proved, sketch]} (rigorous reduction to cited
standard machinery, details routine but not written), or
\emph{[conjecture]}.  ``The 1D Lemma'' means the Lemma of the companion
note: for the escape diffusion
$\dd X=v(X)\circ(\beta\,\dd t+\varepsilon\,\dd W)$,
$v=\sqrt{2(E-V)}$, on $[x_{0},\infty)$ (reflected at $x_{0}$, minimal
at $+\infty$), the boundary $+\infty$ is \emph{regular} iff the
classical time of flight $T=\int^{\infty}\dd x/v$ is finite (minimal
semigroup sub-stochastic and explosive, mean lifetime
$T^{2}$ for $\beta=0$, $T+O(\varepsilon^{2})$ for $\beta=1$) and
\emph{natural} iff $T=\infty$ (conservative and the unique Feller
realization); the driftless speed measure is the WKB deficiency
density $2[2(E-V)]^{-1/2}\dd x$; only continuity of $V$ and eventual
$V<E$ are used.

%======================================================================
\section{What replaces the single end and the single integral}
\label{sec:setting}

In one dimension the escaping end is a point and one integral $T(E)$
decides all three legs.  In the paper's multidimensional setting the
transverse dynamics is resolved into channels: for the separable model
of the channel-additivity theorem (paper Thm.~8.2, \texttt{thm:channels}),
\[
H \;=\; h_{1}\otimes\mathbf 1+\mathbf 1\otimes h_{\perp}
\quad\text{on } L^{2}(\R\times\R^{m}),
\qquad h_{1}=\tfrac12p^{2}+V_{1},\quad
h_{\perp}\varphi_{n}=E_{n}\varphi_{n},
\]
the classical escape mechanism conserves the transverse energy, and the
quantum deficiency vectors are $u_{n}^{\pm}\otimes\varphi_{n}$.  The
correct multidimensional replacements, established below, are:
\begin{itemize}
\item the single escaping end $+\infty$ becomes the \emph{countable set
of channel ends} $\{\infty_{n}\}_{n\in\N}$ (one per transverse
eigenstate); for continuous transverse spectrum, a measure space of
fiber ends (Sec.~\ref{sec:tail});
\item the time of flight becomes the \emph{per-channel} family
\[
T_{n}\;=\;T(E-E_{n})
\;=\;\int_{x_{n}}^{\infty}\frac{\dd x}{\sqrt{2\,(E-E_{n}-V_{1}(x))}},
\qquad
x_{n}\ \text{any base point with }\sup_{x\ge x_{n}}V_{1}<E-E_{n};
\]
this is the precise form taken by the anticipated
``direction-dependent $T(E,\omega)$'': the direction label $\omega$ is
quantized into the channel label $n$ (for central potentials,
literally the angular decomposition, Corollary~\ref{cor:central});
\item and---genuinely new in $d>1$---there is a \emph{second} kind of
end, the \emph{channel tail} $n\to\infty$ at finite or co-moving $x$,
which has no 1D counterpart and which is exactly where the clean
dictionary stops being a theorem (Secs.~\ref{sec:obstruction}--\ref{sec:tail}).
\end{itemize}

Throughout, $V_{1}$ is continuous with $V_{1}(x)\to-\infty$; where the
quantum leg is invoked we assume in addition the paper's Theorem~6.1
regularity.  A first elementary but structurally important fact:

\begin{lemma}[channel class equality; proved]\label{lem:classeq}
If $V_{1}(x)\to-\infty$ then for all $n,m$:
$T_{n}<\infty\iff T_{m}<\infty$.  Hence the \emph{qualitative}
classification is channel-uniform; only the \emph{quantitative} data
$(T_{n})_{n}$ is channel-dependent.
\end{lemma}

\begin{proof}
$\dfrac{E-E_{n}-V_{1}(x)}{E-E_{m}-V_{1}(x)}\to1$ as $x\to\infty$, so the
integrands are comparable near $+\infty$; the finite parts are finite by
continuity and the choice of base points.
\end{proof}

This is the stochastic shadow of the $\lambda$-independence of the
limit-point/limit-circle dichotomy used in the paper's channel-additivity
proof: there, every channel is limit circle because one channel is;
here, every channel end is regular because one is.

%======================================================================
\section{The separable case: full per-channel generalization}
\label{sec:separable}

\paragraph{The channel escape process.}
Define the state space as the disjoint union of half-lines
\[
\calE\;=\;\bigsqcup_{n\in\N} I_{n},\qquad I_{n}=[x_{n},\infty),
\]
and on each $I_{n}$ the channel escape diffusion $X^{(n)}$ of the 1D
Lemma with speed
$v_{n}(x)=\sqrt{2(E-E_{n}-V_{1}(x))}$, clock
$\tau_{n}(x)=\int_{x_{n}}^{x}\dd y/v_{n}(y)$, reflection at $x_{n}$,
minimal (absorbing) convention at $\infty_{n}$; generator
$\mathcal L=\bigoplus_{n}L_{n}$,
$L_{n}=\tfrac{\varepsilon^{2}}{2}(v_{n}\partial_{x})^{2}
+\beta\,v_{n}\partial_{x}$, $\beta\in\{0,1\}$.  This is the stochastic
realization of the separable mechanism: the transverse motion is
bounded oscillation whose only dynamical residue is its conserved
energy, i.e.\ the frozen channel label.

\begin{theorem}[channel dictionary, separable case; proved except (iv)
countable part]\label{thm:A}
Let $V_{1}$ be continuous, $V_{1}\to-\infty$, and let
$(E_{n},\varphi_{n})$ be as above.  Then:
\begin{itemize}
\item[(i)] \textbf{(per-channel classification)}  Each channel end
$\infty_{n}$ is \emph{regular} iff $T_{n}<\infty$ and \emph{natural}
iff $T_{n}=\infty$ (never exit, never entrance); by
Lemma~\ref{lem:classeq} all channel ends are in the same class.  For
$\beta=0$ the Feller test integrals in channel $n$ are
$\Sigma_{n}=N_{n}=T_{n}^{2}$.
\item[(ii)] \textbf{(explosion)}  If $T_{n}<\infty$ (equivalently, for
all $n$): started in channel $n$, the process explodes a.s.\ with mean
lifetime $T_{n}^{2}$ ($\beta=0$) resp.\
$T_{n}-\tfrac{\varepsilon^{2}}{2}(1-e^{-2T_{n}/\varepsilon^{2}})
\in(T_{n}-\tfrac{\varepsilon^{2}}{2},T_{n}]$ ($\beta=1$), and
$\zeta\to T_{n}$ a.s.\ as $\varepsilon\downarrow0$: \emph{each channel
loses its mass at its own classical escape time}.  The minimal
semigroup $\bigoplus_{n}P^{\min,n}_{t}$ is strictly sub-stochastic on
every channel.  If $T_{n}=\infty$ for one (hence every) $n$, the direct
sum is conservative and is the unique Feller realization of
$\mathcal L$ with the given reflecting data.
\item[(iii)] \textbf{(deficiency $=$ speed-measure data, channel by
channel)}  The driftless speed measure of channel $n$ is the WKB
density of the paper's deficiency vector $u_{n}\otimes\varphi_{n}$:
\[
m_{n}(\dd x)\;=\;2\bigl[2(E-E_{n}-V_{1}(x))\bigr]^{-1/2}\dd x
\;\approx\;2\,|u_{n}^{\rm WKB}(x)|^{2}\,\dd x,
\qquad
m_{n}\bigl((x,\infty)\bigr)\approx
2\!\int_{x}^{\infty}\!|u_{n}|^{2},
\]
and the scale mass of channel $n$ near $\infty_{n}$ is $T_{n}$.  Thus
deficiency indices $(\infty,\infty)$ correspond exactly to
\emph{countably many boundary ends, each carrying finite scale and
speed mass}---one Feller datum per deficiency channel.
\item[(iv)] \textbf{(lateral conditions $=$ the $S$-matrix at
infinity)}  Suppose $T_{n}<\infty$ for all $n$, and truncate to $M$
channels.  The sub-Markov Feller extensions of the minimal process are
exactly the Feller--Wentzell lateral conditions at the $M$ regular ends:
at each $\infty_{n}$ a convex datum (absorption, killing, reflection,
stickiness, and a jump-in measure $\nu_{n}$ on $\calE$), the jump-in
measures being allowed to charge \emph{other} channels.  The
conservative extensions are those with no absorption/killing; a packet
escaping in channel $n$ then re-enters distributed according to
$\nu_{n}$---a channel-\emph{mixing} return law, the stochastic mirror
of the paper's unitary $S$-matrix at infinity, and different lateral
data reproduce the regulator-dependence of the returned mixture seen in
the paper's four-channel figure.  \emph{[Finite $M$: proved via
Feller's lateral conditions and the Freidlin--Wentzell gluing theory;
countable $M$: proved-sketch via It\^o excursion theory, with a
summability condition on the excursion entrance laws; the
\emph{completeness} claim (that all Feller extensions arise this way)
is standard for finite $M$ and labeled conjectural for countable
$M$.]}
\end{itemize}
\end{theorem}

\begin{proof}
(i)--(iii) are channelwise reductions.  The direct sum has no
interaction between channels, so every path quantity is computed in the
channel where the process starts; the 1D Lemma applied with potential
$V_{1}+E_{n}$ (equivalently spectral shift $E\mapsto E-E_{n}$) gives
(i) and (ii) verbatim, and Lemma~\ref{lem:classeq} gives the class
equality.  Uniqueness in (ii): at a natural end no lateral condition
exists (Feller), channel by channel, and a Feller extension of a direct
sum that respects the generator must respect the direct-sum structure
away from the boundary; since no boundary is accessible, the minimal
realization is the only one.  For (iii), the deficiency equation in
channel $n$ is $(h_{1}-(\pm\ii-E_{n}))u=0$; to leading WKB order at
real reference energy the modulus is
$|u_{n}|^{2}\approx[2(E-E_{n}-V_{1})]^{-1/2}$ (the shift of the
spectral parameter off the real axis affects the phase, not the leading
modulus---the same convention already used in the paper's one-line
identity), which is $\tfrac12 m_{n}$; the scale mass is
$\int^{\infty}\dd\tau_{n}=T_{n}$ as in 1D.

(iv), finite $M$.  Each $\infty_{n}$ is regular by (i).  Realize
$\calE_{M}=\bigsqcup_{n\le M}I_{n}$, compactify each half-line by its
end, and consider the star obtained by the disjoint boundary set
$\{\infty_{1},\dots,\infty_{M}\}$.  Feller's classification of lateral
conditions at a regular boundary of a generalized second-order operator
[Feller, Illinois J.\ Math.\ 1 (1957) 459--504] provides, for a single
interval, the general boundary datum including a jump-in measure on the
interior; the multi-interval (graph) case, with boundary points allowed
to communicate, is the Freidlin--Wentzell gluing theory [Freidlin and
Wentzell, Ann.\ Probab.\ 21 (1993) 2215--2245], which parametrizes all
Feller processes agreeing with the minimal one in the interior by
Wentzell data at the vertices; jump-in to interior points is supplied by
the It\^o excursion construction of instantaneous-return processes
[It\^o, Proc.\ Sixth Berkeley Symp.\ (1972) 225--239; Blumenthal,
\emph{Excursions of Markov Processes}, Birkh\"auser 1992]: at a regular
end with finite scale and speed mass the minimal process killed at the
end has an exit system, and any sub-probability kernel
$(\nu_{n})_{n\le M}$ together with nonnegative sticky/killing weights
determines a unique right process extending it.  Countable $M$: the
same excursion construction applies end by end; what must be checked is
that the concatenation over countably many ends does not accumulate
infinitely many boundary visits in finite time, for which a uniform
lower bound on mean return times (automatic when
$\inf_{n} T_{n}>0$, and in general a summability condition on the
entrance laws) suffices.  We have not written this in full generality;
the parametrization is asserted as complete only for finite $M$.
\end{proof}

\begin{remark}[the mirror is structural, not a bijection; honest]
The quantum completions form the unitary group
$U(\mathcal K_{+},\mathcal K_{-})$, $\dim\mathcal K_{\pm}=\infty$; the
conservative stochastic completions form the convex set of Markov
gluing data $(\nu_{n})$.  As already in 1D ($U(1)$ versus Feller's
convex family), these are \emph{parallel infinite-dimensional families
attached to the same boundary objects}, in canonical correspondence at
the level of ``which channels communicate with which through
infinity'', not isomorphic as spaces.  In particular the paper's
statement that conservative completions ``scramble the transverse
state with data injected at infinity'' has an exact stochastic
counterpart: the returned channel mixture is the jump-in kernel, and it
is extension data, not data of $\mathcal L$.
\end{remark}

\begin{remark}[conservative / transient / explosive trichotomy]
The paper's Sec.~5 hierarchy is reproduced by the clock picture: in
channel $n$ the drifted model is Brownian motion with unit drift on
$[0,T_{n})$; $T_{n}=\infty$ with $\beta=1$ is \emph{conservative and
transient} (the quadratic-ghost cascade regime, paper
Thm.~\texttt{thm:quadratic}), $T_{n}<\infty$ is \emph{explosive}, and a
confining channel (no escape interval) is \emph{recurrent}.  All three
currencies appear channelwise.
\end{remark}

\begin{remark}[two-ended channels]
If $h_{1}$ is limit circle at both ends (the paper's
\texttt{thm:channels} hypothesis), each channel contributes two regular
ends and the finite-$M$ boundary set has $2M$ points---the exact mirror
of the $(2M,2M)$ indices in the paper's finite-channel stability
theorem (\texttt{thm:stability}).  All statements above apply verbatim
with the reflecting base replaced by the second end's data.
\end{remark}

%======================================================================
\section{Central potentials: the sphere at infinity, resolved into
channels}

\begin{corollary}[central case; proved]\label{cor:central}
Let $H=-\tfrac12\Delta-c|x|^{\alpha}$ on $\R^{d}$, $d\ge2$, $c>0$,
$\alpha>2$ (the paper's Corollary~8.3).  In the angular decomposition
the radial channel $\ell$ carries the escape diffusion with
$v_{\ell}(r)^{2}=2\bigl(E+cr^{\alpha}-L^{2}/(2r^{2})\bigr)$,
$L^{2}=(\ell+\tfrac{d-2}{2})^{2}-\tfrac14$, on $r\ge r_{\ell}$ beyond
the outer turning point.  Then every channel end is regular, with
\[
T_{\ell}\;\asymp\;C_{\alpha}\,c^{-\frac{2}{\alpha+2}}\,
L^{\frac{2-\alpha}{2+\alpha}}\;\xrightarrow[\ell\to\infty]{}\;0,
\qquad
C_{\alpha}=\int_{\rho_{0}}^{\infty}
\bigl(2(\rho^{\alpha}-\tfrac12\rho^{-2})\bigr)^{-1/2}\dd\rho<\infty,
\]
so $T^{*}=\sup_{\ell}T_{\ell}<\infty$: the classification is uniform
across the entire sphere of directions.  All conclusions of
Theorem~\ref{thm:A} hold with $n$ replaced by $\ell$ (finite
multiplicities absorbed into the channel count); in particular the
conservative completions of the central malicious ghost are
angular-momentum-mixing return laws at $r=\infty$.
\end{corollary}

\begin{proof}
Substituting $r=\lambda\rho$ with $\lambda=(L^{2}/c)^{1/(\alpha+2)}$
equalizes the potential and centrifugal scales,
$cr^{\alpha}=c\lambda^{\alpha}\rho^{\alpha}$,
$L^{2}/(2r^{2})=c\lambda^{\alpha}\rho^{-2}/2$, giving
$T_{\ell}=\lambda(2c\lambda^{\alpha})^{-1/2}
\bigl(C_{\alpha}+o(1)\bigr)
=(2c)^{-1/2}\lambda^{1-\alpha/2}(C_{\alpha}+o(1))$ and
$\lambda^{1-\alpha/2}=(L^{2}/c)^{\frac{2-\alpha}{2(\alpha+2)}}$.
(Numerical check, $\alpha=4$, $c=1$: $T_{\ell}$ at
$L^{2}=10^{2},10^{4},10^{6}$ is $1.194,\,0.554,\,0.257$, successive
ratios $2.155$ against the predicted $10^{1/3}=2.154$.)  The
$\rho$-integral converges at infinity because $\alpha>2$ and at the
turning point because the turning point of
$\rho^{\alpha}-\tfrac12\rho^{-2}$ is simple.  The rest is
Theorem~\ref{thm:A}.
\end{proof}

Thus for central potentials the anticipated ``per-direction
$T(E,\omega)$'' collapses, correctly, to the channel family
$T_{\ell}$: the boundary at infinity is not the geometric sphere
$S^{d-1}$ but its harmonic resolution, and the uniformity
$T^{*}<\infty$ is automatic.  This is the multidimensional situation in
which the generalization is \emph{complete}.

%======================================================================
\section{Bounded channel mixing: the stochastic Kato theorem}
\label{sec:bounded}

The paper's non-separable theorem (\texttt{thm:kato}) says a bounded
mixing $B=\mu f(x_{1})g(x_{2})$ preserves the deficiency indices.  The
stochastic counterpart replaces the frozen channel label by a jump
process with \emph{bounded intensity}.

\paragraph{The switched escape process.}  Fix measurable rates
$\lambda(x;n\to m)\ge0$ with
$\sup_{x,n}\sum_{m}\lambda(x;n\to m)\le\kappa<\infty$, supported on
\emph{open} channels ($\lambda(x;n\to m)=0$ unless $x\ge x_{m}$).
Realize the pair $(X_{t},n_{t})$ on a probability space carrying the
driving Brownian motion $W$ and an independent Poisson process of
intensity $\kappa$ with i.i.d.\ uniform marks: between candidate times
the path solves the frozen channel SDE
$\dd X=v_{n}(X)\circ(\beta\,\dd t+\varepsilon\,\dd W)$ (reflected at
$x_{n}$); at a candidate time a switch $n\to m$ is accepted with
probability $\lambda(X;n\to m)/\kappa$ (thinning).  $X$ is continuous
across switches.

\begin{theorem}[stability under bounded mixing; proved]\label{thm:B}
For the switched escape process:
\begin{itemize}
\item[(a)] \textbf{(conservative side, no extra hypothesis)}  If
$T_{n}=\infty$ for one (hence every) channel, the switched process is
conservative: $\Prob(\zeta<\infty)=0$ for every initial condition and
\emph{every} admissible bounded-rate switching law.
\item[(b)] \textbf{(explosive side, uniform time of flight)}  If
\[
(\mathrm U):\qquad T^{*}\;:=\;\sup_{n}T_{n}\;<\;\infty ,
\]
then the switched process explodes a.s.\ for every admissible
bounded-rate law, with exponential tail: for $\beta=1$,
\[
\Prob(\zeta>kh)\le(1-p_{0})^{k},\qquad
h=T^{*}+1,\quad
p_{0}=e^{-\kappa h}\,
\Prob\Bigl(\varepsilon\sup_{s\le h}|W_{s}|\le\tfrac12\Bigr)>0,
\]
so $\E[\zeta]\le h/p_{0}<\infty$; the same holds for $\beta=0$ with
$p_{0}=e^{-\kappa h}\Prob(\sup_{s\le h}W_{s}\ge T^{*})$.
\end{itemize}
Hence under (U) the dictionary survives bounded mixing verbatim:
$T<\infty$ class $\Leftrightarrow$ explosive, $T=\infty$ class
$\Leftrightarrow$ conservative---the stochastic mirror of Kato
stability of the deficiency indices.
\end{theorem}

\begin{proof}
(a)  The candidate switch times form a Poisson($\kappa$) process
independent of everything else, so a.s.\ only finitely many candidates
(hence switches) occur in $[0,t]$.  If $\zeta\le t$, the explosion
instant lies in some inter-candidate interval $[\sigma_{k},
\sigma_{k+1})$; on that interval the path is the frozen channel-$n_k$
diffusion started at the finite point $X_{\sigma_{k}}$, whose end
$\infty_{n_{k}}$ is natural (1D Lemma with $T_{n_{k}}=\infty$), hence
inaccessible in finite time---a null event.  A countable union over $k$
of null events is null.

(b)  Work on windows $[kh,(k+1)h)$.  Let $A_{k}$ be the event
that no Poisson candidate falls in the window and
$\varepsilon\sup_{s\le h}|W_{kh+s}-W_{kh}|\le\tfrac12$;
by independence of the Poisson clock and the Brownian increments from
$\mathcal F_{kh}$, $\Prob(A_{k}\mid\mathcal F_{kh})\ge p_{0}$.  On
$A_{k}$, if $\zeta>kh$, the process sits in some channel $n$ at clock
position $u=\tau_{n}(X_{kh})\in[0,T_{n})$ and, with no switch in the
window, follows the frozen dynamics; in the clock coordinate
(1D Lemma proof) the Skorokhod decomposition gives
$Z_{s}\ge u+s-\tfrac12\ge s-\tfrac12$, and
$s=T_{n}+\tfrac12\le T^{*}+\tfrac12<h$ already forces $Z$ to reach
$T_{n}$: explosion occurs inside the window.  Hence
$\Prob(\zeta>(k+1)h\mid\zeta>kh)\le1-p_{0}$ and the geometric bound
follows; positivity of $p_{0}$ is elementary.  For $\beta=0$ the frozen
clock process is reflected Brownian motion, and
$\sup_{s\le h}W_{s}\ge T^{*}$ forces the hit; the reflection only
helps.  Note the argument never needs to know \emph{which} channels are
visited---only (U) and the open-channel support of the rates (which
guarantees the frozen dynamics is defined after every switch).
\end{proof}

\begin{lemma}[when (U) is automatic; proved]\label{lem:U}
Suppose, for $x\ge R$: $-V_{1}(x)\ge c_{1}x^{2+\epsilon}$ and
$-V_{1}'(x)\ge c_{3}x^{1+\epsilon}$ with $\epsilon>0$, $V_{1}\in C^{1}$
strictly decreasing there, and $T(E)<\infty$.  Then
$\sup_{E'\le E}T(E')<\infty$; in particular (U) holds for every choice
of transverse spectrum $(E_{n})$ bounded below or not.  This covers all
power-law potentials $-V_{1}\sim cx^{2+\epsilon}$, hence the paper's
representative models.
\end{lemma}

\begin{proof}
Write $\mu=-E'\ (\to\infty)$, turning point $x'(E')$, and
$x_{g}=(\mu/c_{1})^{1/(2+\epsilon)}\ge x'(E')$ (since
$E'+c_{1}x'^{2+\epsilon}\le E'-V_{1}(x')=0$).  Beyond $x_{g}$,
$E'-V_{1}\ge c_{1}x^{2+\epsilon}-\mu$ and the scaling
$x=x_{g}u$ gives
$\int_{x_{g}}^{\infty}\le C\,\mu^{-\epsilon/(2(2+\epsilon))}\to0$
exactly as in the central computation.  Between $x'$ and $x_{g}$, the
derivative bound gives
$E'-V_{1}(x)=\int_{x'}^{x}(-V_{1}')\ge c_{3}x'^{1+\epsilon}(x-x')$,
so
$\int_{x'}^{x_{g}}\le
\sqrt{2(x_{g}-x')/(c_{3}x'^{1+\epsilon})}
\le C'\,x_{g}^{-\epsilon/2}\to0$
(using $x'\asymp x_{g}$).  So $T(E')\to0$ as $E'\to-\infty$, and on any
compact range of $E'$, $T$ is finite and continuous (simple turning
points by strict monotonicity, dominated convergence with the local
$(x-x')^{-1/2}$ bound).  Hence the sup is finite.
\end{proof}

%======================================================================
\section{The obstruction: unbounded channel transfer
(``channel riding'')}
\label{sec:obstruction}

Everything so far generalizes.  What follows is the exact point of
failure, and it is proved by construction, not merely feared.

\begin{lemma}[channel-riding reduction; proved]\label{lem:riding}
Let $n(\cdot):\bigl[x_{*},\infty\bigr)\to\N$ be measurable and locally
constant off a discrete set, with channel $n(x)$ open at $x$ for every
$x$.  The \emph{position-locked} switched process (switch to channel
$n(x)$ upon reaching $x$; a limit of bounded-rate laws with rates
$\to\infty$, so \emph{outside} Theorem~\ref{thm:B}'s hypothesis) is the
one-dimensional generalized diffusion on $[x_{*},\infty)$ with
effective speed
\[
v_{\rm eff}(x)\;=\;\sqrt{2\bigl(E-E_{n(x)}-V_{1}(x)\bigr)},
\]
defined by scale $s_{\rm eff}(x)=\int^{x}\dd y/v_{\rm eff}$ (for
$\beta=0$; the drifted variant analogously) and speed measure
$m_{\rm eff}=2\,\dd x/v_{\rm eff}$.  Its end at $+\infty$ is classified
by the 1D Lemma with the \emph{effective time of flight}
\[
T_{\rm eff}\;=\;\int_{x_{*}}^{\infty}
\frac{\dd x}{\sqrt{2\,(E-E_{n(x)}-V_{1}(x))}}\,.
\]
\end{lemma}

\begin{proof}
Scale-and-speed constructions require only a strictly increasing
continuous scale and a Radon speed measure (It\^o--McKean \S4);
$v_{\rm eff}$ is measurable, locally bounded and locally bounded away
from $0$ on the open region, so $s_{\rm eff}$ and $m_{\rm eff}$
qualify, and the resulting regular diffusion coincides with the
switched process between switch points by locality of the generator.
The classification integrals are computed exactly as in the 1D Lemma
($\Sigma=N=T_{\rm eff}^{2}$ for $\beta=0$).
\end{proof}

$T_{\rm eff}$ is a \emph{path functional of the coupled mechanism}, not
a function of the per-channel family $\{T_{n}\}$.  The next two
propositions make the underdetermination rigorous, in both directions.

\begin{proposition}[Counterexample I: all channels explosive, coupled
process conservative; proved]\label{prop:ce1}
There exist $V_{1}$ (continuous, strictly decreasing to $-\infty$,
satisfying the paper's Theorem~6.1 regularity, including
$|V_{1}'||V_{1}|^{-3/2}\to0$) and a transverse spectrum
$(E_{n})$, such that
\[
\sup_{n}T_{n}\;\le\;\tfrac32
\qquad(\text{every frozen channel explosive, uniformly---(U) holds}),
\]
yet the position-locked riding law $n(x)$ below gives
$T_{\rm eff}=\infty$: the coupled process is \emph{conservative}.
Consequently the per-channel dictionary
(``$T_{n}<\infty\ \forall n\Rightarrow$ explosive'') is \emph{false}
for unbounded-rate channel transfer, even under (U); the bounded-rate
hypothesis of Theorem~\ref{thm:B} cannot be dropped.
\end{proposition}

\begin{proof}
Take $E=0$, $E_{n}=16^{n}$, and the shelf potential
\[
V_{1}(x)=-E_{n}-\tfrac12 \ \ \text{on } [2n,\,2n+1],\qquad
V_{1} \text{ linear on } [2n+1,\,2n+2],\qquad n\ge1,
\]
(extended arbitrarily and boundedly to the left; tilt the shelves by an
immaterial slope if strict monotonicity is wanted).  Regularity:
$|V_{1}'|\asymp16^{n}$ on the $n$-th drop while
$|V_{1}|^{3/2}\asymp16^{3n/2}$, so
$|V_{1}'||V_{1}|^{-3/2}\to0$.  Channel $m$ is open exactly from its own
shelf onward ($E-E_{m}-V_{1}=E_{n}+\tfrac12-E_{m}>0$ iff $n\ge m$),
so $x_{m}=2m$, and
\[
T_{m}\;\le\;
\underbrace{\frac{1}{\sqrt{2\cdot\tfrac12}}}_{\text{own shelf}}
+\underbrace{1}_{\text{own drop, speed}^{2}\ge1}
+\sum_{n>m}\Bigl(\underbrace{\tfrac{1}{\sqrt{2(E_{n}-E_{m})}}}_{\text{shelf }n}
+\underbrace{\tfrac{2}{\sqrt{E_{n}}}}_{\text{drop }n}\Bigr)
\;\le\;\tfrac32
\]
for $m$ large (numerically $T_{1}=1.17$, $T_{3}=1.01$,
$T_{8}=1.00004$; all $\le1.18$).  Now ride: $n(x)=n$ on
$[2n,2n+2)$ (always open).  On shelf $n$ the effective speed is
$\sqrt{2(E_{n}+\tfrac12-E_{n})}=1$ over unit length, so each shelf
contributes exactly $1$ to $T_{\rm eff}$:
$T_{\rm eff}\ge\sum_{n}1=\infty$.  Lemma~\ref{lem:riding} and the 1D
Lemma ($T_{\rm eff}=\infty$ $\Rightarrow$ natural, conservative)
finish.
\end{proof}

\begin{proposition}[Counterexample II: all channels conservative,
coupled process explosive---ghost-specific; proved]\label{prop:ce2}
Let $V_{1}\equiv0$ on $[1,\infty)$ (or $V_{1}=-x^{2}$, the marginal
case), $E=1$, and let $h_{\perp}$ be a \emph{ghost} oscillator,
$E_{n}=-(n+\tfrac12)$, unbounded below.  Every frozen channel has
constant (or at-most-linearly-growing) speed, hence
$T_{n}=\infty$ for all $n$: every frozen channel is conservative.  The
riding law $n(x)=\lceil x^{2+\epsilon}\rceil$ gives
$v_{\rm eff}(x)^{2}=2(E+n(x)+\tfrac12)\ge2x^{2+\epsilon}$, so
$T_{\rm eff}\le\int^{\infty}x^{-1-\epsilon/2}\dd x/\sqrt2<\infty$:
the coupled process is \emph{explosive}.  The mechanism extracts
unbounded energy from the transverse ghost---the stochastic realization
of the paper's ``channel-tail escape'' ($\lambda<0$ PU climbing a
rising potential through arbitrarily high mode-2 channels) and of its
remark that transverse observables are not $H_{0}$-bounded for
indefinite $h_{\perp}$.
\end{proposition}

\begin{proof}
Immediate from Lemma~\ref{lem:riding} and the 1D Lemma.
\end{proof}

\begin{theorem}[ghosts are necessary for Counterexample II; proved
under (D)]\label{thm:F}
Suppose $h_{\perp}$ is bounded below, $E_{n}\ge E_{0}$, let
$\bar v(x)^{2}=2(E-E_{0}-V_{1}(x))$ dominate every channel speed, and
assume
\[
(\mathrm D):\qquad \int^{\infty}\frac{\dd x}{\bar v}=\infty
\quad\text{and}\quad
\limsup_{x\to\infty}\bar v'(x)<\infty
\ \ \bigl(\text{i.e.\ }-V_{1}'\le C\,\bar v \text{ eventually};
\text{ holds for }V_{1}\ge-Kx^{2}\bigr).
\]
Then \emph{every} switched escape process---arbitrary adapted switching
with countably many switches, including position-locked riding and any
unbounded rates---is conservative.  In particular, below-quadratic
(marginal or slower) escape ends with non-ghost transverse dynamics are
\emph{absolutely} stable under channel coupling: the conservative leg
of the dictionary is robust, the explosive leg is not.
\end{theorem}

\begin{proof}
Khasminskii test with the dominating clock.  Let
$\bar\tau(x)=\int_{x_{*}}^{x}\dd y/\bar v$ and let $\varphi$ equal
$\bar\tau$ for $x\ge x_{*}+1$, be smoothed to constant near the base so
that $\varphi'\ge0$ vanishes at every reflection point;
$\varphi\uparrow\infty$ by (D).  For any channel $n$ open at $x$
(so $v_{n}(x)\le\bar v(x)$), on $x\ge x_{*}+1$:
$v_{n}\varphi'=v_{n}/\bar v\le1$, and, using
$v_{n}v_{n}'=-V_{1}'=\bar v\,\bar v'$ (all channels share
$-V_{1}'$),
\[
v_{n}\,( v_{n}\varphi')'
= \frac{v_{n}v_{n}'}{\bar v}-\frac{v_{n}^{2}\,\bar v'}{\bar v^{2}}
= \bar v'\Bigl(1-\frac{v_{n}^{2}}{\bar v^{2}}\Bigr)
\;\in\;\bigl[\min(0,\bar v'),\ \max(0,\bar v')\bigr]
\;\le\;C .
\]
Hence $L_{n}\varphi\le\beta+\tfrac{\varepsilon^{2}}{2}C=:C'$ uniformly
in $n$ and $x$ (on the compact modification region all terms are
bounded since $v_{n}\le\bar v$ there too).  $X$ is continuous across
switches and $\varphi'$ vanishes at reflection points, so
$\varphi(X_{t\wedge\zeta_{k}})-C't$ is a supermartingale for the
switched process up to the exit time $\zeta_{k}$ of $[x_{*},k]$;
$\E\varphi(X_{t\wedge\zeta_{k}})\le\varphi(x_{0})+C't$ while
$\varphi(k)\to\infty$ gives
$\Prob(\zeta_{k}\le t)\le(\varphi(x_{0})+C't)/\varphi(k)\to0$: no
explosion.
\end{proof}

\begin{remark}[exactly where separability is and is not needed;
summary of Secs.~2--5]
\emph{Not needed}: separability is not needed for the per-channel
statements to survive bounded mixing (Theorem~\ref{thm:B}), nor for
the conservative side at all when $h_{\perp}$ is bounded below and the
end is at most quadratic (Theorem~\ref{thm:F}: no hypothesis on the
coupling whatsoever).  \emph{Needed}: some control of channel transfer
(bounded rate, or transverse boundedness below) is needed for the
explosive side; Proposition~\ref{prop:ce1} shows the per-channel family
$\{T_{n}\}$---even uniformly bounded---does not determine the coupled
classification once transfer is unbounded, and
Proposition~\ref{prop:ce2} shows the conservative side also fails, but
\emph{only} by extracting unbounded transverse energy, i.e.\ only for
ghosts (Theorem~\ref{thm:F}).  The deciding object in the coupled
setting is the time-of-flight functional $T_{\rm eff}$ \emph{along the
actual coupled escape mechanism}---one integral per escape
\emph{path family}, no longer one integral per channel.
\end{remark}

%======================================================================
\section{The second boundary: the channel tail, and the exact open
point}
\label{sec:tail}

The riding constructions show that in $d>1$ the state space
$\{(x,n)\}$ has, besides the countable set of channel ends at
$x=\infty$, a second end: $n\to\infty$ (at co-moving $x$).  It is
itself a Feller boundary with a one-sum criterion.

\begin{proposition}[tail end of the channel chain; proved
(classical)]\label{prop:tail}
Model the transverse transfer alone as a minimal birth chain on $\N$
with upward rates $q_{n}>0$.  The tail $n=\infty$ is reached in finite
time a.s.\ iff $\sum_{n}1/q_{n}<\infty$, and the mean climbing time is
exactly $\sum_{n}1/q_{n}$ (Feller's classical explosion criterion for
birth processes).  Thus the tail end carries its own ``time of
flight'' $T_{\rm tail}=\sum_{n}1/q_{n}$, in exact formal analogy with
$T=\int\dd x/v$ at a spatial end---\emph{the dictionary acquires a
second leg of the same shape}, but its datum $q_{n}$ comes from the
coupling matrix elements, not from any integral of the potential.
\end{proposition}

\begin{conjecture}[coupled dictionary; clean statement of what remains]
\label{conj:coupled}
For the switched escape process with adapted rates
$\lambda(x;n\to m)$ derived from a (possibly unbounded) coupling, the
process explodes iff the combined time-of-flight functional is finite
along some accessible escape route:
\[
\Prob(\zeta<\infty)>0
\iff
\inf_{\text{admissible routes }(x(t),n(t))}
\Bigl[\,T_{\rm eff}\bigl(x(\cdot),n(\cdot)\bigr)\Bigr]<\infty ,
\]
where routes carry both spatial ends (classified by $T_{\rm eff}$,
Lemma~\ref{lem:riding}) and the tail end (classified by
$T_{\rm tail}$-type sums, Proposition~\ref{prop:tail}).  All results
above are the proved instances: separable (routes are single channels),
bounded rate (routes cannot ride), riding (routes computed
explicitly).  The general statement is a conjecture because
accessibility of a route for the \emph{given} rates is a
large-deviations question we have not settled.
\end{conjecture}

\begin{remark}[what this says about the paper's open problem; honest]
The paper's remaining gap---deficiency of the full quartic PU, i.e.\
whether the finite-channel indices $(2M,2M)$ ($\lambda>0$) resp.\
$(0,0)$ ($\lambda<0$) survive $M\to\infty$
(\texttt{thm:stability}, \texttt{cor:dichotomy},
\texttt{prop:directintegral})---is the exact quantum image of the
dichotomy proved here:
\begin{itemize}
\item $\lambda>0$ PU: classical escape locks onto a definite kinetic
direction (paper: form ratio $\to0.9997$), i.e.\ bounded channel
occupation along the escape; that is the regime of
Theorem~\ref{thm:B}(b)/(U), where explosivity (deficiency) is stable.
The stochastic analysis therefore \emph{supports} deficiency
$(\infty,\infty)$ for $\lambda>0$, but does not prove it: the
interchannel couplings are polynomial in $k$, not bounded, so
Theorem~\ref{thm:B} does not apply literally---the needed strengthening
(subordinate-but-unbounded transfer with bounded occupation along
escape) is exactly the open tail control.
\item $\lambda<0$ PU: the classical escape is channel-tail riding
(paper: climbs through arbitrarily high mode-2 channels), the regime of
Proposition~\ref{prop:ce2}/Theorem~\ref{thm:F}; there the classical and
quantum legs \emph{can} and apparently \emph{do} split (classical
blowup, quantum essential self-adjointness at every finite $M$, drift
evidence for the full operator): zero-point transverse energy
renormalizes the effective rates $q_{n}$, closing the tail.  Our
framework locates the split precisely at the tail end and explains its
mechanism (the tail leg's datum is coupling-derived and
quantization-sensitive, unlike the potential-derived $T_{n}$); it does
not decide the full-operator question, which remains open, as in the
paper.
\end{itemize}
Note the direction of robustness proved here (conservative side stable
under arbitrary coupling for non-ghost transverse dynamics,
Theorem~\ref{thm:F}; explosive side fragile, Proposition~\ref{prop:ce1})
matches the asymmetry of the paper's numerical verdicts: the
completion candidates are both on the ``classically explosive but
quantum-safe'' side, never the reverse.
\end{remark}

\begin{remark}[continuous transverse spectrum / sphere at infinity;
conjecture]
If $h_{\perp}$ has continuous spectrum (the direct-integral situation
of the paper's \texttt{prop:directintegral}), the channel ends form a
measure space of fibers.  Per-fiber classification holds a.e.\
(Lemma~\ref{lem:classeq} and the 1D Lemma fiberwise, proved), and
measurable fields of lateral conditions mirror measurable fields of
deficiency solutions; but individual fiber ends have zero capacity for
the coupled process, so lateral data must be specified as kernels up to
null sets, and the transverse operator in the fiber frame is unbounded
(the paper's proposition)---the coupled classification is again
governed by tail control, and we do not claim it.  This, not the
geometry of $S^{d-1}$, is the honest form of the ``sphere at
infinity'': a direct integral of clocks, with the same open tail.
\end{remark}

%======================================================================
\section*{Verdict}

\textbf{PER-CHANNEL-ONLY} (with sharp boundaries in both directions):

\begin{itemize}
\item \textbf{Proved, full generalization where the paper is
separable/central/boundedly mixed.}  The escape-diffusion dictionary
extends verbatim channel by channel: the escaping end becomes the
countable set of channel ends, $T(E)$ becomes $T_{n}=T(E-E_{n})$, each
end is regular iff $T_{n}<\infty$ (all in the same class), per-channel
mean explosion time $T_{n}^{2}$ resp.\ $T_{n}+O(\varepsilon^{2})$,
per-channel speed measure $=$ the WKB density of the paper's deficiency
vector $u_{n}\otimes\varphi_{n}$, and deficiency $(\infty,\infty)$
corresponds to the infinite family of Feller--Wentzell lateral
conditions with channel-mixing jump-in kernels---the stochastic
$S$-matrix at infinity (Theorem~\ref{thm:A}; Corollary~\ref{cor:central}
with uniformity $T_{\ell}\propto\ell^{(2-\alpha)/(2+\alpha)}$;
Theorem~\ref{thm:B} $=$ the stochastic Kato theorem, under the uniform
time-of-flight condition (U), automatic for power-law ends,
Lemma~\ref{lem:U}).
\item \textbf{Proved obstruction.}  ``One integral per channel'' is the
\emph{most} that survives: the family $\{T_{n}\}$ does not determine
the coupled classification under unbounded channel transfer.
Counterexample I (shelf potential, $\sup_n T_n\le\frac32$, riding gives
$T_{\rm eff}=\infty$): all channels explosive, coupled process
conservative.  Counterexample II (ghost transverse spectrum): all
channels conservative, coupled process explosive---and this direction
is provably \emph{ghost-specific} (Theorem~\ref{thm:F}: impossible when
$h_{\perp}$ is bounded below and the end is at most quadratic).  The
deciding object in genuine multi-D is the time-of-flight functional
along the coupled escape route, plus a second, tail boundary at
$n=\infty$ with its own one-sum criterion
(Proposition~\ref{prop:tail}).
\item \textbf{Open / conjectural, matching the paper's open problem.}
The route-infimum dictionary (Conjecture~\ref{conj:coupled}); the
countable-boundary completeness of the lateral-condition
parametrization; and the quantum tail (full PU at $M\to\infty$), where
the classical and quantum legs provably can split and the stochastic
leg must be run with quantization-corrected transfer rates.
\end{itemize}

\begin{thebibliography}{9}
\bibitem{Feller1957} W.~Feller, \emph{Generalized second order
differential operators and their lateral conditions}, Illinois J.\
Math.\ \textbf{1}, 459--504 (1957).
\bibitem{ItoMcKean} K.~It\^o and H.~P.~McKean, Jr., \emph{Diffusion
Processes and their Sample Paths}, Springer (1974), \S\S3.7--3.8,
4.1--4.6.
\bibitem{FW1993} M.~I.~Freidlin and A.~D.~Wentzell, \emph{Diffusion
processes on graphs and the averaging principle}, Ann.\ Probab.\
\textbf{21}, 2215--2245 (1993).
\bibitem{Ito1972} K.~It\^o, \emph{Poisson point processes attached to
Markov processes}, Proc.\ Sixth Berkeley Symp.\ Math.\ Statist.\
Prob.\ III, 225--239 (1972); R.~M.~Blumenthal, \emph{Excursions of
Markov Processes}, Birkh\"auser (1992).
\bibitem{KarlinTaylor} S.~Karlin and H.~M.~Taylor, \emph{A Second
Course in Stochastic Processes}, Academic Press (1981), Ch.~15
(boundary classification; explosion criteria; birth-process explosion,
\S4).
\bibitem{Khas} R.~Z.~Khasminskii, \emph{Stochastic Stability of
Differential Equations}, 2nd ed., Springer (2012), Thm.~3.5
(non-explosion Lyapunov test).
\bibitem{ReedSimon2} M.~Reed and B.~Simon, \emph{Methods of Modern
Mathematical Physics II}, Academic Press (1975), Thms.~X.7--X.10.
\bibitem{KatoBook} T.~Kato, \emph{Perturbation Theory for Linear
Operators}, Springer (1980), Thm.~V.4.3 (stability of deficiency
indices).
\end{thebibliography}

\end{document}
